\chapter{Proposed Architecture for a Digital Twin-Based eHealth System}

\section{Architecture Rationale}
Based on the literature review, a Digital Twin-based eHealth system should not be limited to a data dashboard or a predictive algorithm. It should be organized as a layered system that connects the physical healthcare environment to a virtual representation, processes heterogeneous data, provides simulation and analysis services, and ensures security, privacy, and interoperability.

The proposed architecture is general and can be adapted to different eHealth scenarios such as remote monitoring, chronic disease management, hospital workflow optimization, rehabilitation support, wellness tracking, or clinical decision support. It is not designed for one specific disease or one specific project.

\section{Overview of the Proposed Architecture}
The proposed architecture is composed of seven layers: physical layer, data acquisition layer, interoperability and preprocessing layer, data management layer, virtual twin layer, intelligence and simulation layer, and service layer. Security, privacy, and governance are considered as cross-cutting concerns across all layers.

\begin{figure}[H]
    \centering
    \resizebox{0.98\textwidth}{!}{\begin{tikzpicture}[font=\small, >=Latex]
\tikzstyle{layer}=[rounded corners=6pt, draw=DTGray, thick, fill=DTLightGray, minimum width=13.2cm, minimum height=0.95cm, align=center]
\tikzstyle{security}=[rounded corners=8pt, draw=DTOrange, thick, dashed, fill=DTLightOrange, inner sep=0.25cm]

\node[layer, fill=DTLightBlue, draw=DTBlue] (service) at (0,0) {\textbf{Service Layer}: dashboards, alerts, reports, decision support, scenario simulation};
\node[layer, fill=DTLightPurple, draw=DTPurple] (ai) at (0,-1.25) {\textbf{Intelligence and Simulation Layer}: AI/ML, rules, prediction, optimization, uncertainty};
\node[layer, fill=DTLightGreen, draw=DTGreen] (virtual) at (0,-2.5) {\textbf{Virtual Twin Layer}: dynamic virtual representation, state update, context model};
\node[layer] (data) at (0,-3.75) {\textbf{Data Management Layer}: databases, time-series storage, metadata, audit logs};
\node[layer] (interop) at (0,-5.0) {\textbf{Interoperability and Preprocessing}: cleaning, normalization, HL7/FHIR mapping};
\node[layer] (acq) at (0,-6.25) {\textbf{Data Acquisition Layer}: connected devices, mobile apps, EHR systems, manual inputs};
\node[layer, fill=DTLightOrange, draw=DTOrange] (physical) at (0,-7.5) {\textbf{Physical Layer}: patients, clinicians, medical devices, healthcare environments};

\draw[->, thick, DTBlue] (physical) -- (acq);
\draw[->, thick, DTBlue] (acq) -- (interop);
\draw[->, thick, DTBlue] (interop) -- (data);
\draw[->, thick, DTBlue] (data) -- (virtual);
\draw[->, thick, DTBlue] (virtual) -- (ai);
\draw[->, thick, DTBlue] (ai) -- (service);
\draw[->, thick, DTGreen, dashed] (service.east) .. controls (7.2,-1.0) and (7.2,-6.5) .. node[right, align=center]{feedback / recommendations} (physical.east);

\begin{scope}[on background layer]
\node[security, fit=(service)(ai)(virtual)(data)(interop)(acq)(physical), label={[text=DTOrange]right:\rotatebox{-90}{\textbf{Security, privacy, governance, validation}}}] {};
\end{scope}
\end{tikzpicture}
}
    \caption{Proposed multi-layer architecture for a Digital Twin-based eHealth system. Author synthesis based on the reviewed literature.}
    \label{fig:proposed_architecture}
\end{figure}

\section{Physical Layer}
The physical layer includes real-world entities involved in the eHealth ecosystem. These may include patients, healthcare professionals, medical devices, wearable sensors, mobile phones, hospital equipment, and healthcare environments. The role of this layer is to generate real-world data and receive possible feedback or recommendations from the digital services.

\section{Data Acquisition Layer}
The data acquisition layer collects information from connected devices, mobile applications, clinical systems, and manual inputs. Data may be continuous, periodic, or event-based. Examples include physiological measurements, activity data, medication information, symptoms, clinical notes, laboratory results, and operational data.

\section{Interoperability and Preprocessing Layer}
This layer prepares data for use by the Digital Twin. It handles data cleaning, normalization, timestamp alignment, missing value management, anomaly detection, and format conversion. It also supports interoperability standards such as HL7 and FHIR to connect the Digital Twin with healthcare information systems.

\section{Data Management Layer}
The data management layer stores historical, real-time, and contextual information. It may include relational databases, time-series databases, clinical repositories, data warehouses, and metadata catalogs. Data governance rules must define who can access data, how long data are stored, and how data are anonymized or pseudonymized when needed.

\section{Virtual Twin Layer}
The virtual twin layer represents the current state of the physical entity or healthcare process. Depending on the application, the virtual twin may represent a patient profile, a disease trajectory, a medical device, a care pathway, a hospital service, or a public health process. The virtual twin should be continuously updated when new valid data become available.

\section{Intelligence and Simulation Layer}
This layer provides analytical and predictive capabilities. It may include rule-based reasoning, machine learning models, deep learning models, statistical models, mechanistic simulations, optimization algorithms, and uncertainty estimation. The purpose is to transform the virtual representation into useful knowledge for monitoring, prediction, scenario analysis, and decision support.

\section{Service Layer}
The service layer exposes the functions of the Digital Twin to users and systems. Possible services include dashboards, alerts, reports, risk assessment, scenario simulation, remote monitoring, workflow optimization, and decision support. These services must be designed according to user needs, clinical safety, explainability, and usability requirements.

\section{Security, Privacy, and Governance Layer}
Security and privacy are cross-cutting concerns. The system should include authentication, authorization, encryption, audit logs, secure APIs, role-based access control, data minimization, consent management, and governance policies. Ethical oversight is also required to ensure fairness, transparency, and responsible use of the Digital Twin.

\begin{figure}[H]
    \centering
    \resizebox{0.9\textwidth}{!}{\begin{tikzpicture}[font=\small, >=Latex, node distance=1.0cm]
\tikzstyle{center}=[circle, draw=DTBlue, very thick, fill=DTLightBlue, minimum size=2.8cm, align=center]
\tikzstyle{box}=[rounded corners=6pt, draw=DTGray, thick, fill=DTLightGray, minimum width=3.1cm, minimum height=0.9cm, align=center]
\tikzstyle{arrow}=[->, thick, DTGray]

\node[center] (dt) {eHealth\\\textbf{Digital Twin}};
\node[box, above=1.4cm of dt] (auth) {Authentication\\authorization};
\node[box, below=1.4cm of dt] (privacy) {Data minimization\\consent management};
\node[box, left=3.6cm of dt] (encrypt) {Encryption\\secure APIs};
\node[box, right=3.6cm of dt] (audit) {Audit logs\\traceability};
\node[box, above left=1.1cm and 2.5cm of dt] (govern) {Governance\\policies};
\node[box, above right=1.1cm and 2.5cm of dt] (ethics) {Ethics\\fairness, transparency};
\node[box, below left=1.1cm and 2.5cm of dt] (valid) {Validation\\clinical plausibility};
\node[box, below right=1.1cm and 2.5cm of dt] (monitor) {Monitoring\\incident response};

\foreach \x in {auth,privacy,encrypt,audit,govern,ethics,valid,monitor}
  \draw[arrow] (\x) -- (dt);
\end{tikzpicture}
}
    \caption{Security, privacy, and governance concerns around the eHealth Digital Twin. Author synthesis based on the reviewed literature.}
    \label{fig:security_privacy}
\end{figure}

\section{Data Flow}
Figure~\ref{fig:data_flow} illustrates the general data flow. Real-world data are collected from the physical layer, transmitted through acquisition modules, preprocessed and standardized, stored in the data layer, and used to update the virtual twin. The intelligence and simulation layer generates insights that are exposed through services. Feedback may then be provided to patients, clinicians, or healthcare systems depending on the use case.

\begin{figure}[H]
    \centering
    \resizebox{0.98\textwidth}{!}{\begin{tikzpicture}[font=\small, >=Latex, node distance=1.2cm]
\tikzstyle{box}=[rounded corners=6pt, draw=DTGray, thick, fill=DTLightGray, minimum width=2.7cm, minimum height=1.0cm, align=center]
\tikzstyle{bluebox}=[rounded corners=6pt, draw=DTBlue, thick, fill=DTLightBlue, minimum width=2.7cm, minimum height=1.0cm, align=center]
\tikzstyle{greenbox}=[rounded corners=6pt, draw=DTGreen, thick, fill=DTLightGreen, minimum width=2.7cm, minimum height=1.0cm, align=center]
\tikzstyle{purplebox}=[rounded corners=6pt, draw=DTPurple, thick, fill=DTLightPurple, minimum width=2.7cm, minimum height=1.0cm, align=center]

\node[box] (patient) {Physical world\\patient / process};
\node[bluebox, right=1.0cm of patient] (sources) {Data sources\\devices, apps, EHR};
\node[bluebox, right=1.0cm of sources] (preprocess) {Preprocessing\\clean, align, validate};
\node[box, right=1.0cm of preprocess] (storage) {Data storage\\history, context};
\node[greenbox, below=1.5cm of storage] (twin) {Virtual Twin\\state and behavior};
\node[purplebox, left=1.0cm of twin] (analytics) {Analytics\\AI + simulation};
\node[purplebox, left=1.0cm of analytics] (services) {Services\\dashboard, alerts};
\node[box, left=1.0cm of services] (users) {Users\\clinicians, patients};

\draw[->, thick] (patient) -- (sources);
\draw[->, thick] (sources) -- (preprocess);
\draw[->, thick] (preprocess) -- (storage);
\draw[->, thick] (storage) -- (twin);
\draw[->, thick] (twin) -- (analytics);
\draw[->, thick] (analytics) -- (services);
\draw[->, thick] (services) -- (users);
\draw[->, thick, dashed, DTGreen] (users.north) .. controls (-3.5,2.0) and (-5.2,1.8) .. node[above, align=center]{feedback / actions} (patient.north);
\draw[->, thick, dashed, DTBlue] (analytics.north) -- node[right]{update} (storage.south);
\end{tikzpicture}
}
    \caption{General data flow in a Digital Twin-based eHealth system. Author synthesis based on the reviewed literature.}
    \label{fig:data_flow}
\end{figure}

\section{General Use Case Scenario}
A general eHealth use case can be described as follows. A patient uses connected devices and a mobile application to collect health-related data at home. These data are transmitted to the eHealth platform, cleaned, standardized, and stored. The virtual twin is updated with the new measurements and historical context. The intelligence layer analyzes trends, detects anomalies, estimates risk, and simulates possible scenarios. The service layer displays a dashboard, generates alerts when necessary, and produces a report for healthcare professionals.

This scenario remains general and can be adapted to different domains such as chronic disease follow-up, elderly care, rehabilitation, mental health monitoring, wellness tracking, hospital-at-home programs, and clinical workflow optimization.

\section{Evaluation Plan}
The proposed architecture can be evaluated using several criteria:

\begin{itemize}[leftmargin=*]
    \item \textbf{Functional evaluation:} ability to collect, store, update, analyze, and display data.
    \item \textbf{Interoperability evaluation:} compatibility with standards and existing healthcare systems.
    \item \textbf{Performance evaluation:} latency, scalability, and reliability.
    \item \textbf{Security evaluation:} authentication, authorization, encryption, and auditability.
    \item \textbf{Usability evaluation:} clarity of dashboards and usefulness for users.
    \item \textbf{Clinical plausibility:} relevance and safety of generated outputs.
    \item \textbf{Ethical evaluation:} fairness, transparency, consent, and data governance.
\end{itemize}

\section{Conclusion}
This chapter proposed a general conceptual architecture for Digital Twin-based eHealth systems. The architecture is designed to be reusable, modular, secure, interoperable, and adaptable to multiple healthcare scenarios.

\clearpage
